% ООП ФИИТ: Информационные технологии в управлении и принятии решений
% (направление 02.04.02 «Фундаментальная информатика и информационные технологии»)
\program{mitcd}

\title{%
Определение структурных сегментов и эмоциональных характеристик
музыкальных композиций методами машинного обучения%
\english
Structural Segmentation and Emotional Characterisation of Musical
Compositions by Machine Learning Methods%
}

\student{%
Антонов Михаил Евгеньевич%
\brief М.\,Е.\,Антонов%
\group М-26%
\english Antonov Mikhail Evgenyevich%
}

\superviser{%
Кудряшов Максим Юрьевич%
\brief М.\,Ю.\,Кудряшов%
\position доцент кафедры информационных технологий%
\degree кандидат физико-математических наук%
}

\year{2025-2026}

\goal{%
Разработка и экспериментальное исследование программного модуля
автоматического анализа музыкальных композиций, выполняющего
структурную сегментацию, оценку эмоциональных характеристик и
жанровую классификацию аудиозаписей на основе мультизадачной
модели глубокого обучения.%
}

% ---------------------------------------------------------------------------
% Задачи: перечень переписывается в подраздел «Цель и задачи» и повторно
% используется в подразделе «Основные результаты». Формат элемента:
%   \task{<формулировка задачи>\result{<краткий результат>}}
% ---------------------------------------------------------------------------
\begin{tasks}
\task{%
Провести обзор области Music Information Retrieval и обосновать
мультизадачный подход к объединению четырёх задач.%
\result{%
Выполнен обзор 40 источников по трём задачам MIR
и пяти семействам аудио-архитектур (Глава~1). По итогам
обзора в качестве методологической основы выбрана
мультизадачная модель с общим аудио-энкодером и четырьмя
независимыми классификационными головами. Аргументы в её
пользу (содержательная связь задач, регуляризующий эффект
общего ствола и экономия вычислений при инференсе) собраны
в~\textsection\,2.3.%
}}

\task{%
Подготовить обучающий корпус на основе трёх открытых датасетов
(DEAM, Harmonix~Set, GTZAN) с унифицированным пайплайном
извлечения признаков.%
\result{%
Собран объединённый корпус из 3\,545 композиций на основе
трёх открытых датасетов: DEAM (эмоции), Harmonix~Set
(структура, с автоматическим скачиванием 743~треков
с YouTube через \texttt{yt-dlp}) и GTZAN (жанр). Реализован
модуль \texttt{musicml.features} для унифицированного
извлечения логарифмических мел-спектрограмм со 128~фильтрами,
нарезки на десятисекундные окна с шагом 1~секунда и
глобальной нормализации. Итоговый обучающий набор насчитывает
176\,127 окон, разбиение на обучающую, валидационную и
тестовую выборки 70/15/15 делается на уровне треков, что
исключает утечку между сплитами.%
}}

\task{%
Разработать и реализовать пять архитектур мультизадачных
моделей: CNN, CNN\,+\,BiLSTM, PANNs\,+\,Linear,
PANNs\,+\,BiLSTM и AST.%
\result{%
Реализованы и обучены пять архитектур, охватывающих ключевые
оси сравнения. В их числе: свёрточная сеть с блоками
канального внимания (около 800~тысяч параметров), её надстройка
двунаправленной LSTM, линейный классификатор и BiLSTM поверх
замороженных эмбеддингов PANNs CNN14, а также частично
размороженный Audio Spectrogram Transformer
с 14{,}4\,млн обучаемых параметров из 87\,млн общих.
У всех конфигураций единый интерфейс выхода: четыре
классификационные головы плюс две регрессионные для непрерывных
значений возбуждения и валентности. Обучение всех пяти
проведено на одном тестовом корпусе с единым протоколом,
чтобы результаты были напрямую сопоставимы.%
}}

\task{%
Выполнить экспериментальное сравнение пяти архитектур по четырём
задачам на едином тестовом корпусе.%
\result{%
Проведены четыре последовательных эксперимента: базовый
запуск v1, улучшенный v2, расширение трансформер-архитектурой
AST~v2 и неудачный CNN~v3 с агрессивным перевзвешиванием
классов (документированный negative result). Лучший по
среднему результату оказался AST~v2: точность 58{,}5\,\%
усреднённо по четырём задачам, 40{,}6\,\% на структурной
сегментации (на 5{,}3 процентных пункта выше CNN~v2) и
81{,}5\,\% на жанровой классификации. Победители по отдельным
задачам распределились между AST~v2 (сегментация),
PANNs+BiLSTM (возбуждение и валентность) и CNN~v2 (жанр).
Одной архитектуры, доминирующей сразу по всем четырём задачам,
не нашлось.%
}}

\task{%
Реализовать постобработку сегментных предсказаний на базе
декодера Витерби с музыкологическими априорными распределениями.%
\result{%
Разработан и реализован декодер~v2 для постобработки сегментных
предсказаний. В нём объединены шесть механизмов: температурное
масштабирование логитов, асимметричная матрица переходов
6$\times$6 с музыкологическими априори, позиционные приоры
для вступления и концовки, поклассовые ограничения минимальной
длительности, привязка границ к локальным максимумам кривой
новизны и проверка повторяющихся сегментов на близость
представлений. Полный пайплайн повышает Boundary~F1 с 0{,}31
до 0{,}45 (на 14 процентных пунктов выше базового аргмакса)
и сокращает среднее число выдаваемых сегментов
с 28 до 7--8 на трёхминутный трек.%
}}

\task{%
Разработать веб-прототип для интерактивной демонстрации
результатов.%
\result{%
Реализован полнофункциональный веб-прототип с трёхзвенной
архитектурой: фронтенд на React\,19 с Vite, бэкенд на Elysia
поверх Bun, сервис инференса на FastAPI с предобученной AST~v2
в памяти. Центральное решение интерфейса~--- синхронизированный
дашборд из шести панелей с общим временным курсором, в который
сведены разметка структуры, эмоциональные кривые, жанровое
распределение, мел-спектрограмма и сводная панель характеристик.
Дополнительно реализована автоматическая подгрузка
метаданных треков из ID3-тегов и YouTube.%
}}

\task{%
Сформулировать ограничения подхода и направления дальнейшей
работы.%
\result{%
Систематически проанализированы пять групп ограничений текущего
подхода. К ним относятся смещение жанровой головы на современной
музыке из-за устаревшей таксономии GTZAN, фундаментальный потолок
шестиклассовой сегментации в районе 50\,\% по причинам данных
и разметки, размер AST в 574\,МБ с восьмисекундным временем
инференса на трек, невозможность мульти-жанровых предсказаний
в текущей постановке и затруднения с очень длинными треками.
Для каждой группы предложены конкретные направления развития:
надстройка sequence-level декодера BiLSTM или CRF поверх
AST-эмбеддингов, дообучение на современных датасетах
FMA-medium или SALAMI, дистилляция знаний из AST в более
компактный CNN, потоковый инференс со скользящим окном.%
}}
\end{tasks}
